\documentclass{article}

\begin{document}
    \begin{center}
         \Large \textbf{Challenge 6: Custom Images} \\ [6pt]
        \normalsize
        Autor: Maximilian Jakob Maag B.Sc. \\
        Version: 1.0
    \end{center}

    \section{Intro}
    Immer wieder wird unter dem Betriff Golden Image oder Custom Image verlangt ein geändertes Linux Image zu deployen.
    In dieser Challenge können Sie das nötige Vorgehen üben.

    \section{Lernziele}
    Sie lernen ein eignes minimales Linux ISO zu erstellen.
    Diese Fähigkeit soll Ihnen helfen custom images bereitstellen zu können.

    \section{Aufgaben}
    Diese Aufgaben führen Sie zum Deployment eines kleinen custom images.

    \subsection{Aufgabe 1}
    Untersuchen Sie die Größe eines typischen Ubuntu Images. Warum ist das ISO File so groß?

    \subsection{Aufgabe 2}
    Wie groß ist der Linux Kernel selbst? Warum unterscheiden sich die Angabe, die SIe dazu finden können?

    \subsection{Aufgabe 3}
    Nehmen Sie ein Ubuntu Image. Verändern Sie dieses Minimal, packen Sie es in ein neues ISO file und starten Sie eine VM damit.
    Verifizieren Sie Ihre Änderung, indem Sie die Hashwerte Ihres Images mit dem Ausgangs Image vergleichen.

    \subsection{Aufgabe 4}
    Erstellen Sie ein custom Image mit Linux Kernel und Shell das nicht größer als 1 GiB ist.

    \subsection{Aufgabe 5}
    Erstellen Sie das kleinstmögliche deploybare Linux Image. Es muss den Kernel, eine Shell und ein Anwendungsprogramm enthalten.
    Das Image sollte eine Größe von 2 MB nicht übersteigen. TIPP: Es ist möglich ein Linuxsystem von einer Floppidisk to booten.
    
    \subsection{Aufgabe 6}
    Erstellen Sie ein Linux Image das nach dem Booten nur lesenden Zugriff auf das Filesystem hat.
    Welche Anwendungen könnten hier Realisiert werden? TIPP: Read Only Linuxsysteme kommen oft in Firewallsystemen zum Einsatz.
    Entwicklen Sie einen Usecase, erstellen Sie ein passendes Image und deployen Sie es auf geeigneter Infrastruktur.
\end{document}